% !TEX program = xelatex
\documentclass[11pt,a4paper]{ctexart}

% ==================== 颜色 ====================
\usepackage{xcolor}
\definecolor{maincolor}{HTML}{1A7A6B}       % 主色（标题、顶栏、框边框）
\definecolor{accent}{HTML}{D4A843}          % 强调色（关键词高亮）
\definecolor{boxbg}{HTML}{F0F7F5}           % 问题框背景
\definecolor{borderblue}{HTML}{2E86AB}      % 方法框边框
\definecolor{proofbg}{HTML}{FDF8F0}         % 方法框背景
\definecolor{summarybg}{HTML}{F5F0FF}       % 总结框背景

% ==================== 页面 ====================
\usepackage[top=2.2cm, bottom=2.2cm, left=2.5cm, right=2.5cm, headheight=23.12pt]{geometry}

% ==================== 字体 ====================
\setmainfont{TeX Gyre Pagella}

% ==================== 数学 ====================
\usepackage{amsmath,amssymb,amsfonts,mathtools}

% ==================== 彩色框 ====================
\usepackage{tcolorbox}
\tcbuselibrary{skins,breakable,most}

% ==================== 页眉页脚 ====================
\usepackage{fancyhdr}
\pagestyle{fancy}
\fancyhf{}
\fancyhead[L]{\color{maincolor}\sffamily\bfseries 数学讲义}
\fancyhead[R]{\color{gray}\sffamily \today}
\fancyfoot[C]{\color{gray}\sffamily\thepage}
\renewcommand{\headrulewidth}{1.5pt}
\renewcommand{\headrule}{\hbox to\headwidth{\color{maincolor}\leaders\hrule height \headrulewidth\hfill}}

% ==================== 章节标题样式 ====================
\usepackage{titlesec}
\titleformat{\section}
  {\color{maincolor}\sffamily\LARGE\bfseries}
  {\thesection}{1em}{}
  [\vspace{0.3cm}{\color{maincolor!40}\titlerule[2pt]}]

\titleformat{\subsection}
  {\color{maincolor!80}\sffamily\Large\bfseries}
  {\thesubsection}{1em}{}

% ==================== 表格 ====================
\usepackage{array,booktabs,colortbl}

% ==================== 杂项 ====================
\usepackage{enumitem}
\setlist{nosep, leftmargin=*}
\usepackage[hidelinks]{hyperref}

% ==================== 自定义命令 ====================
\newcommand{\highlight}[1]{\colorbox{accent!20}{\sffamily\bfseries #1}}

% ==================== 自定义环境 ====================

% 1. 问题框 —— 青绿边框 + 浅绿背景
\newtcolorbox[auto counter]{problembox}[1][]{
  enhanced,
  colframe=maincolor,
  colback=boxbg,
  coltitle=white,
  fonttitle=\sffamily\bfseries,
  title={\large 问题},
  attach boxed title to top left={xshift=4mm, yshift=-2mm},
  boxed title style={colback=maincolor, sharp corners},
  sharp corners,
  boxrule=1.2pt,
  left=4mm, right=4mm, top=5mm, bottom=4mm,
  before skip=1.2em, after skip=1.2em,
  #1
}

% 2. 方法/证明框 —— 蓝边框 + 暖白背景
\newtcolorbox[auto counter]{methodbox}[2][]{
  enhanced,
  colframe=borderblue,
  colback=proofbg,
  coltitle=white,
  fonttitle=\sffamily\bfseries,
  title={#2},
  attach boxed title to top left={xshift=4mm, yshift=-2mm},
  boxed title style={colback=borderblue, sharp corners},
  sharp corners,
  boxrule=1.2pt,
  left=4mm, right=4mm, top=5mm, bottom=4mm,
  before skip=1.2em, after skip=1.2em,
  #1
}

% 3. 总结框 —— 淡紫背景
\newtcolorbox{summarybox}[1][]{
  enhanced,
  colframe=maincolor!60,
  colback=summarybg,
  coltitle=white,
  fonttitle=\sffamily\bfseries,
  title={\large $\blacktriangleright$ 总结},
  attach boxed title to top left={xshift=4mm, yshift=-2mm},
  boxed title style={colback=maincolor!70, sharp corners},
  sharp corners,
  boxrule=1.2pt,
  left=4mm, right=4mm, top=5mm, bottom=4mm,
  before skip=1.5em, after skip=1.2em,
  #1
}

% ==================== 文档开始 ====================
\begin{document}

% ---- 彩色顶栏标题 ----
\begin{tcolorbox}[
  enhanced,
  colframe=maincolor,
  colback=maincolor,
  sharp corners,
  boxrule=0pt,
  height=2.8cm,
  valign=center,
  before skip=-2.2cm,
  after skip=1.5cm,
  grow to left by=2.5cm,
  grow to right by=2.5cm,
]
  \begin{center}
    {\color{white}\sffamily\bfseries\Huge 表达式整数解问题}\\[3pt]
    {\color{white!85}\sffamily\large 求使 $(b-\frac{1}{a})(c-\frac{1}{b})(a-\frac{1}{c})$ 为整数的正整数解}
  \end{center}
\end{tcolorbox}

\vspace{0.5cm}

\begin{flushright}
  {\color{gray}\sffamily \today}
\end{flushright}

% ==================== 正文 ====================

\section{问题陈述}

\begin{problembox}
求所有正整数三元组 $(a,b,c)$，使得表达式
\[
  E = \left(b - \frac{1}{a}\right)\left(c - \frac{1}{b}\right)\left(a - \frac{1}{c}\right)
\]
的值为\highlight{整数}。
\end{problembox}

解题思路：先将表达式展开并分离整数部分与分数部分，再利用对称性和不等式放缩，分情况讨论得到所有解。

\section{展开与分离}

\begin{methodbox}{核心思路}
将原式展开后，显式地分离\highlight{整数部分}与\highlight{分数部分}。因为整个 $E$ 为整数，且整数部分自动为整数，所以分数部分也必须为整数。
\end{methodbox}

\subsection*{展开原式}

将三个括号乘开：
\[
  E = abc - a - b - c + \frac{1}{a} + \frac{1}{b} + \frac{1}{c} - \frac{1}{abc}
\]

\subsection*{分离整数部分与分数部分}

将明显为整数的项归为一组，含分母的项归为另一组：
\[
  E = (abc - a - b - c) + \left(\frac{1}{a} + \frac{1}{b} + \frac{1}{c} - \frac{1}{abc}\right)
\]

令分数部分为 $k$：
\[
  k = \frac{1}{a} + \frac{1}{b} + \frac{1}{c} - \frac{1}{abc}
\]

由于 $abc-a-b-c$ 是整数，$E$ 为整数\highlight{当且仅当 $k$ 为整数}。

\section{利用对称性简化}

\begin{methodbox}{核心思路}
表达式关于 $a,b,c$ 完全对称，因此可以假设 $1 \le a \le b \le c$ 讨论，最后通过全排列恢复所有解。
\end{methodbox}

\subsection*{对称性假设}

由于原表达式关于 $a,b,c$ 对称，不妨假设：
\[
  1 \le a \le b \le c
\]

在此假设下求出所有解后，再对解进行全排列即可得到全部答案。

\section{情况一：$a = 1$}

\begin{methodbox}{核心思路}
将 $a=1$ 代入 $k$ 的表达式，通过不等式放缩确定分式部分只能取 $1$，进而解出 $b$ 和 $c$。
\end{methodbox}

\subsection*{代入 $a=1$}

\[
  k = 1 + \frac{1}{b} + \frac{1}{c} - \frac{1}{bc}
    = 1 + \frac{b + c - 1}{bc}
\]

由于 $b,c \ge 1$，分式部分 $\frac{b+c-1}{bc} > 0$。

\subsection*{确定分式部分的取值}

由恒等变形：
\[
  (b-1)(c-1) \ge 0 \;\Longrightarrow\; bc \ge b + c - 1
\]

因此：
\[
  0 < \frac{b+c-1}{bc} \le 1
\]

既然该分式是整数 $k$ 减去 $1$，且取值严格在 $(0,1]$ 之间，它只能等于 $1$。

\subsection*{求解}

\[
  \frac{b + c - 1}{bc} = 1 \;\Longrightarrow\; b + c - 1 = bc
\]
\[
  bc - b - c + 1 = 0 \;\Longrightarrow\; (b-1)(c-1) = 0
\]

结合 $b \le c$ 且 $b \ge a = 1$，解得：
\[
  b = 1
\]

此时 $c$ 为任意正整数。即得到解族：
\[
  (a,b,c) = (1, 1, m) \quad (m \in \mathbb{Z}^+)
\]

% ==================== 新页 ====================
\newpage

\section{情况二：$a \ge 2$}

\begin{methodbox}{核心思路}
当 $a \ge 2$ 时，通过放缩确定 $k$ 的取值范围，锁定 $k=1$，再将条件转化为关于 $a,b,c$ 的方程，通过凑因式分解求解。
\end{methodbox}

\subsection*{放缩确定 $k$ 的值}

此时有 $2 \le a \le b \le c$，对 $k$ 进行上下界放缩：
\[
  0 < k \le \frac{1}{2} + \frac{1}{2} + \frac{1}{2} - \frac{1}{8}
       = \frac{11}{8} < 2
\]

$k$ 为正整数且严格小于 $2$，故：
\[
  k = 1
\]

\subsection*{转化为整系数方程}

由 $k=1$ 即：
\[
  \frac{1}{a} + \frac{1}{b} + \frac{1}{c} - \frac{1}{abc} = 1
\]

两边同乘 $abc$ 并移项：
\[
  bc + ac + ab - 1 = abc
\]
\[
  \boxed{\,abc - ab - bc - ca = -1\,}
\]

\section{凑因式分解求解}

\begin{methodbox}{核心思路}
取 $a=2$，代入后将方程配方为 $(b-2)(c-2)=3$，利用 $b \le c$ 解出唯一解。再说明 $a \ge 3$ 无解。
\end{methodbox}

\subsection*{代入 $a=2$}

将 $a=2$ 代入 $abc - ab - bc - ca = -1$：
\[
  2bc - 2b - bc - 2c = -1 \;\Longrightarrow\; bc - 2b - 2c = -1
\]

两边同加 $4$ 凑成可因式分解的形式：
\[
  bc - 2b - 2c + 4 = 3 \;\Longrightarrow\; (b-2)(c-2) = 3
\]

\subsection*{求解 $(b-2)(c-2)=3$}

$3$ 的正因子对为 $(1,3)$。结合 $b \le c$：
\[
  \begin{cases}
    b-2 = 1 \\
    c-2 = 3
  \end{cases}
  \;\Longrightarrow\;
  b = 3,\; c = 5
\]

得到一组解 $(a,b,c) = (2, 3, 5)$。

\subsection*{$a \ge 3$ 无解}

同理可证，当 $a \ge 3$ 时方程无正整数解（过程与 $a=2$ 类似，此处从略）。

\section{总结}

\begin{summarybox}

\subsection*{推导结构}

\begin{center}
\begin{tabular}{@{}>{\sffamily\bfseries}p{3.2cm} p{10.5cm}@{}}
\toprule
\textbf{步骤} & \textbf{内容} \\
\midrule
展开分离 & $E = (abc-a-b-c) + \bigl(\frac{1}{a}+\frac{1}{b}+\frac{1}{c}-\frac{1}{abc}\bigr)$，分数部分 $k$ 必须为整数 \\
对称假设 & 不妨设 $1 \le a \le b \le c$，最后全排列恢复 \\
$a=1$ 情况 & $k=1+\frac{b+c-1}{bc}$，由 $bc \ge b+c-1$ 得 $\frac{b+c-1}{bc}=1$，解出 $b=1$ \\
& $\;\to\;$ 解族 $(1,1,m),\ m\in\mathbb{Z}^+$ \\
$a \ge 2$ 情况 & 放缩得 $0<k\le\frac{11}{8}<2$，故 $k=1$，化为 $abc-ab-bc-ca=-1$ \\
代入 $a=2$ & 配方得 $(b-2)(c-2)=3$，解出 $b=3,c=5$ \\
$a \ge 3$ & 无正整数解 \\
全排列 & 消除对称假设，对所有解进行排列 \\
\bottomrule
\end{tabular}
\end{center}

\vspace{1em}

\subsection*{结论}

满足条件的所有正整数解为：

\begin{center}
\begin{tcolorbox}[
  enhanced, colframe=accent!80, colback=accent!10,
  sharp corners, boxrule=1.5pt, width=14cm, center,
]
  \begin{center}
    {\color{maincolor}\sffamily\bfseries\large
      $(a,b,c) = (1, 1, m),\ (1, m, 1),\ (m, 1, 1) \quad (m \in \mathbb{Z}^+)$
    }\\[6pt]
    {\color{maincolor}\sffamily\bfseries\large
      $(a,b,c) = (2, 3, 5),\ (2, 5, 3),\ (3, 2, 5),\ (3, 5, 2),\ (5, 2, 3),\ (5, 3, 2)$
    }
  \end{center}
\end{tcolorbox}
\end{center}

验证：以 $(2,3,5)$ 为例，
\[
  E = \left(3 - \frac{1}{2}\right)\left(5 - \frac{1}{3}\right)\left(2 - \frac{1}{5}\right)
    = \frac{5}{2} \cdot \frac{14}{3} \cdot \frac{9}{5}
    = \frac{5 \cdot 14 \cdot 9}{2 \cdot 3 \cdot 5}
    = \frac{630}{30} = 21 \in \mathbb{Z}
\]

\end{summarybox}

\end{document}
